\documentclass{article}

% Recommended packages for better typesetting:
\usepackage{microtype}
\usepackage{graphicx}
\usepackage{subcaption}
\usepackage{booktabs}

% hyperref makes hyperlinks in the resulting PDF.
% If your build breaks, comment out hyperref below and replace
% \usepackage{icml2026} with \usepackage[nohyperref]{icml2026} above.
\usepackage{hyperref}

% Attempt to make hyperref and algorithmic work together better:
\newcommand{\theHalgorithm}{\arabic{algorithm}}

% Use the following line for the initial blind version submitted for review:
\usepackage{icml2026}

% For preprint, use:
% \usepackage[preprint]{icml2026}

% If accepted, instead use the following line for the camera-ready submission:
% \usepackage[accepted]{icml2026}

\usepackage{amsmath}
\usepackage{amssymb}
\usepackage{mathtools}
\usepackage{amsthm}

% if you use cleveref:
\usepackage[capitalize,noabbrev]{cleveref}

% The \icmltitle you define below is probably too long as a header.
% Therefore, a short form for the running title is supplied here:
\icmltitlerunning{EEG-VAE: Short Running Title}

\begin{document}
\onecolumn

\icmltitle{EEG-VAE: Full Paper Title}

% It is OKAY to include author information, even for blind submissions:
% the style file will automatically remove it unless you've provided
% the [accepted] option to the icml2026 package.

\begin{icmlauthorlist}
  \icmlauthor{Prénom Nom}{inst1}
\end{icmlauthorlist}

\icmlaffiliation{inst1}{Département, Université, Pays}
\icmlcorrespondingauthor{Prénom Nom}{email@exemple.com}

\icmlkeywords{EEG, VAE, generative model, BCI}

\vskip 0.3in

% This command creates the footnote in the first column listing affiliations.
% Use ONE of the following lines. DO NOT remove the command.
\printAffiliationsAndNotice{}  % no special notice (required even if empty)
% \printAffiliationsAndNotice{\icmlEqualContribution}  % for equal contribution

\begin{abstract}
Ton résumé ici (150–200 mots).
\end{abstract}


\section{Introduction}

- EEG

Contributions :
- goal of the paper : propose a variational autoencoder (VAE) for EEG data called EEG-VAE
- propose a 2 stage pretraining pipeline to enforce the learninf of more semantic representations.
- analyze the relevance of the mamba architecture for EEG data.
- evaluate the learned representations with a faire and unbiased comparison with other foundation models with a focus on imagined speech decoding modality.


\section{Method}

\subsection{Pipeline Overview}

\subsection{EEG-VAE Architecture}

\subsection{Losses}

\subsection{Pretraining Results and Analysis}


\section{Downstream Evaluation and Comparisons}


\section{Conclusion}


\newpage
\nocite{*}
\bibliography{references}
\bibliographystyle{icml2026}



\newpage
\appendix
\onecolumn
\section{Foundation models used for comparisons}


\section{Datasets details and preprocessing}

\section{Computational analysis}

\section{Implementation details}


\end{document}
